\chapter{Commit Stage}

\section{Overview}

The Commit stage is the final stage of the Sargantana pipeline and is responsible for updating
the architectural state in program order. Although instructions may execute out of order, they
must always retire sequentially to preserve the behavior defined by the RISC-V ISA.

Only instructions that have successfully completed execution and whose older instructions have
already retired are allowed to commit. This guarantees precise exceptions and ensures that
speculative execution remains invisible to software.

The commit stage therefore represents the point at which speculative results become
architecturally permanent.

\begin{rtlrefs}
  \rtlitem{rtl/datapath/rtl/wb_stage/rtl/graduation_list.sv}
  \rtlitem{rtl/control_unit/rtl/control_unit.sv}
  \rtlitem{rtl/datapath/rtl/datapath.sv}
\end{rtlrefs}

\begin{figure}[h]
  \centering
  \includegraphics[width=\textwidth]{ch12_commit_overview.png}
  \caption{Commit stage in context, highlighting the decode-to-commit control path}
  \label{fig:commit-overview}
\end{figure}

\section{Reorder Buffer (ROB)}

Although the repository does not contain a standalone module named \rtl{rob.sv}, the
reorder/retirement functionality is implemented through the \textbf{Graduation List}, together with the
rename structures and control logic.

The Graduation List maintains the order of all in-flight instructions from dispatch until retirement.

For each instruction it tracks information such as:
\begin{itemize}
  \item destination physical register
  \item completion status
  \item exception status
  \item branch information
  \item retirement eligibility
\end{itemize}

When execution finishes, the corresponding entry is marked as completed. The Commit stage
continuously checks the oldest valid entry. If the instruction has completed without exceptions, it
is retired.

Thus, the Graduation List performs the architectural role traditionally associated with a Reorder
Buffer (ROB).

\begin{rtlrefs}
  \rtlitem{rtl/datapath/rtl/wb_stage/rtl/graduation_list.sv}
  \rtlitem{rtl/datapath/rtl/ir_stage/rtl/rename_table.sv}
  \rtlitem{rtl/datapath/rtl/ir_stage/rtl/free_list.sv}
\end{rtlrefs}

\begin{figure}[h]
  \centering
  \includegraphics[width=\textwidth]{ch12_reorder_buffer.png}
  \caption{Graduation List acting as the Reorder Buffer: recovery and commit signal paths highlighted}
  \label{fig:reorder-buffer}
\end{figure}

\section{In-Order Retirement}

Although execution is highly parallel, retirement always occurs in program order.

For each instruction, the Commit stage performs several checks before retirement:
\begin{itemize}
  \item Has execution completed?
  \item Did the instruction raise an exception?
  \item Is it the oldest instruction?
  \item Are all previous instructions already retired?
\end{itemize}

Only after satisfying these conditions are the architectural structures updated.

Retirement includes:
\begin{itemize}
  \item updating the architectural register mapping
  \item releasing obsolete physical registers
  \item committing memory operations
  \item updating CSR state if required
  \item removing the instruction from the Graduation List
\end{itemize}

This mechanism ensures that software observes execution exactly as if instructions had been
executed sequentially.

\begin{rtlrefs}
  \rtlitem{rtl/datapath/rtl/wb_stage/rtl/graduation_list.sv}
  \rtlitem{rtl/datapath/rtl/ir_stage/rtl/free_list.sv}
  \rtlitem{rtl/datapath/rtl/ir_stage/rtl/rename_table.sv}
\end{rtlrefs}

\section{Exception Handling}

One of the primary responsibilities of the Commit stage is precise exception handling.

If an instruction reports an exception during execution, retirement stops when that instruction
reaches the head of the Graduation List.

The processor then:
\begin{itemize}
  \item flushes younger speculative instructions
  \item restores the correct architectural state
  \item redirects the PC to the exception handler
  \item updates the required CSRs
  \item resumes execution from the exception vector
\end{itemize}

Since no younger instruction has committed yet, software observes a perfectly precise
exception.

This behavior is one of the major advantages of separating execution from retirement.

\begin{rtlrefs}
  \rtlitem{rtl/control_unit/rtl/control_unit.sv}
  \rtlitem{rtl/datapath/rtl/wb_stage/rtl/graduation_list.sv}
  \rtlitem{rtl/datapath/interface_csr/rtl/csr_interface.sv}
\end{rtlrefs}

\section{Pipeline Flush and Recovery}

Whenever speculation proves incorrect, the Commit stage coordinates recovery.

Typical recovery events include:
\begin{itemize}
  \item Branch misprediction
  \item Exceptions
  \item Interrupts
  \item System reset
  \item CSR redirection
\end{itemize}

Recovery generally involves:
\begin{enumerate}
  \item stopping retirement,
  \item flushing younger instructions,
  \item restoring the correct rename state,
  \item recovering the free list,
  \item redirecting the Program Counter,
  \item restarting instruction fetch.
\end{enumerate}

Because retirement is always in-order, recovery restores the processor to a precise architectural
state.

\begin{rtlrefs}
  \rtlitem{rtl/control_unit/rtl/control_unit.sv}
  \rtlitem{rtl/datapath/rtl/if_stage_1/rtl/branch_predictor.sv}
  \rtlitem{rtl/datapath/rtl/if_stage_1/rtl/if_stage_1.sv}
  \rtlitem{rtl/datapath/rtl/wb_stage/rtl/graduation_list.sv}
  \rtlitem{rtl/datapath/rtl/ir_stage/rtl/rename_table.sv}
  \rtlitem{rtl/datapath/rtl/ir_stage/rtl/free_list.sv}
\end{rtlrefs}
